\documentclass[10pt,a4paper]{amsart}
\usepackage{amsmath,amssymb,amsthm,graphicx}
% Retained minimal preamble for the Stacks Project chapter sites-cohomology.tex (Cohomology on Sites).
% Verbatim excerpt of the chapter preamble: Xy-pic setup (native preamble lines 14-18) and the native
% theorem declarations with the mathematical macros (native preamble lines 154-217). Copied mechanically;
% nothing added, removed or redefined. The wrapper class and the amsmath/amssymb/amsthm/graphicx base
% packages are supplied by the assembler, so the class/externaldocument part of the native preamble is omitted.
% For commutative diagrams we use Xy-pic
\usepackage[all]{xy}

% We use 2cell for 2-commutative diagrams.
\xyoption{2cell}
% Theorem environments.
%
\theoremstyle{plain}
\newtheorem{theorem}[subsection]{Theorem}
\newtheorem{proposition}[subsection]{Proposition}
\newtheorem{lemma}[subsection]{Lemma}

\theoremstyle{definition}
\newtheorem{definition}[subsection]{Definition}
\newtheorem{example}[subsection]{Example}
\newtheorem{exercise}[subsection]{Exercise}
\newtheorem{situation}[subsection]{Situation}

\theoremstyle{remark}
\newtheorem{remark}[subsection]{Remark}
\newtheorem{remarks}[subsection]{Remarks}

\numberwithin{equation}{subsection}

% Macros
%
\def\lim{\mathop{\mathrm{lim}}\nolimits}
\def\colim{\mathop{\mathrm{colim}}\nolimits}
\def\Spec{\mathop{\mathrm{Spec}}}
\def\Hom{\mathop{\mathrm{Hom}}\nolimits}
\def\Ext{\mathop{\mathrm{Ext}}\nolimits}
\def\SheafHom{\mathop{\mathcal{H}\!\mathit{om}}\nolimits}
\def\SheafExt{\mathop{\mathcal{E}\!\mathit{xt}}\nolimits}
\def\Sch{\mathit{Sch}}
\def\Mor{\mathop{\mathrm{Mor}}\nolimits}
\def\Ob{\mathop{\mathrm{Ob}}\nolimits}
\def\Sh{\mathop{\mathit{Sh}}\nolimits}
\def\NL{\mathop{N\!L}\nolimits}
\def\CH{\mathop{\mathrm{CH}}\nolimits}
\def\proetale{{pro\text{-}\acute{e}tale}}
\def\etale{{\acute{e}tale}}
\def\QCoh{\mathit{QCoh}}
\def\Ker{\mathop{\mathrm{Ker}}}
\def\Im{\mathop{\mathrm{Im}}}
\def\Coker{\mathop{\mathrm{Coker}}}
\def\Coim{\mathop{\mathrm{Coim}}}

% Boxtimes
%
\DeclareMathSymbol{\boxtimes}{\mathbin}{AMSa}{"02}

%
% Macros for moduli stacks/spaces
%
\def\QCohstack{\mathcal{QC}\!\mathit{oh}}
\def\Cohstack{\mathcal{C}\!\mathit{oh}}
\def\Spacesstack{\mathcal{S}\!\mathit{paces}}
\def\Quotfunctor{\mathrm{Quot}}
\def\Hilbfunctor{\mathrm{Hilb}}
\def\Curvesstack{\mathcal{C}\!\mathit{urves}}
\def\Polarizedstack{\mathcal{P}\!\mathit{olarized}}
\def\Complexesstack{\mathcal{C}\!\mathit{omplexes}}
% \Pic is the operator that assigns to X its picard group, usage \Pic(X)
% \Picardstack_{X/B} denotes the Picard stack of X over B
% \Picardfunctor_{X/B} denotes the Picard functor of X over B
\def\Pic{\mathop{\mathrm{Pic}}\nolimits}
\def\Picardstack{\mathcal{P}\!\mathit{ic}}
\def\Picardfunctor{\mathrm{Pic}}
\def\Deformationcategory{\mathcal{D}\!\mathit{ef}}
\title{On a ringed site the O-modules that are summands of finite free O-modules form a category equivalent to the finite projective R-modules for R the ring of global sections}
\author{The Stacks Project Authors}
\date{stacks-project:sites-cohomology.tex:local-snapshot-20261011; GNU FDL 1.2 or later}
\begin{document}
\maketitle
\noindent\textit{Editorial note.} All delivered mathematical text is the author's own, copied from the pinned Stacks Project chapter file sites-cohomology.tex; no statement or proof step was rewritten, reordered or removed.\par\medskip
\noindent\textit{Editorial note.} The counted proof is the author's complete argument: Same-name restatement check: the label lemma-category-summands-finite-free also occurs in the chapter cohomology.tex (ringed spaces, a special case of the ringed-site setting) and is there a special case, so the more general ringed-site statement is the one counted for this lane.\par\medskip
\noindent\textit{Editorial note.} The label lemma-category-summands-finite-free defines the statement; the counted proof is the proof environment immediately following the lemma in the source.\par\medskip
\noindent\textit{Editorial note.} The preamble is the chapter's own retained minimal preamble (sources/stacks-sites-cohomology/retained-preamble.tex), a verbatim excerpt of the native preamble: the Xy-pic setup and the native theorem declarations with the mathematical macros.\par\medskip
\noindent\textit{Allowed standard prerequisites (editorial).}\begin{itemize}
\item Global sections functor on a ringed site. Same chapter, the section functor Gamma(C, -) on O-modules of a ringed site (C, O) with values in R-modules, R = Gamma(C, O). Supplies the functor E -> Gamma(C, E) and its essential surjectivity target.
\item Constant sheaf of a set on a point topos and the morphism of ringed topoi. Same chapter, the pushforward f\_* = global sections and the constant-sheaf pullback f\textasciicircum{}\{-1\} defining the morphism of ringed topoi (Sh(C), O) -> (Sh(pt), R). Supplies the inverse functor f\textasciicircum{}* = Gamma(C, -) tensored with O, used to reconstruct a module from its finite projective sections.
\end{itemize}\medskip
\begin{lemma}
\label{lemma-category-summands-finite-free}
Let $(\mathcal{C}, \mathcal{O})$ be a ringed site.
Set $R = \Gamma(\mathcal{C}, \mathcal{O})$. The category of
$\mathcal{O}$-modules which are summands of finite free
$\mathcal{O}$-modules is equivalent to the category of
finite projective $R$-modules.
\end{lemma}

\begin{proof}
Observe that a finite projective $R$-module is the same thing
as a summand of a finite free $R$-module.
The equivalence is given by the functor $\mathcal{E} \mapsto
\Gamma(\mathcal{C}, \mathcal{E})$.
The inverse functor is given by the following construction.
Consider the morphism of topoi $f : \Sh(\mathcal{C}) \to \Sh(\text{pt})$
with $f_*$ given by taking global sections and
$f^{-1}$ by sending a set $S$, i.e., an object of
$\Sh(\text{pt})$, to the constant sheaf with value $S$.
We obtain a morphism
$(f, f^\sharp) : (\Sh(\mathcal{C}), \mathcal{O}) \to (\Sh(\text{pt}), R)$
of ringed topoi by using the identity map $R \to f_*\mathcal{O}$.
Then the inverse functor is given by $f^*$.
\end{proof}

\end{document}
